\chapter{In Which Monsieur Fouquet Acts}

\lettrine{I}{n} the meantime Fouquet was hastening to the Louvre, at  the best speed of his English horses. The king was at work with Colbert. All at  once the king became thoughtful. The two sentences of death he had signed on  mounting his throne sometimes recurred to his memory; they were two black spots  which he saw with his eyes open; two spots of blood which he saw when his eyes  were closed. <Monsieur,> said he rather sharply, to the intendant;  <it sometimes seems to me that those two men you made me condemn were not  very great culprits.>

<Sire, they were picked out from the herd of the farmers of the  financiers, which wanted decimating.>

<Picked out by whom?>

<By necessity, sire,> replied Colbert, coldly.

<Necessity!—a great word,> murmured the young king.

<A great goddess, sire.>

<They were devoted friends of the superintendent, were they not?>

<Yes, sire; friends who would have given up their lives for Monsieur  Fouquet.>

<They have given them, monsieur,> said the king.

<That is true;—but uselessly, by good luck,—which was not  their intention.>

<How much money had these men fraudulently obtained?>

<Ten millions, perhaps; of which six have been confiscated.>

<And is that money in my coffers?> said the king with a certain air  of repugnance.

<It is there, sire; but this confiscation, whilst threatening M. Fouquet,  has not touched him.>

<You conclude, then, M. Colbert\longdash>

<That if M. Fouquet has raised against your majesty a troop of factious  rioters to extricate his friends from punishment, he will raise an army when he  has in turn to extricate himself from punishment.>

The king darted at his confidant one of those looks which resemble the livid  fire of a flash of lightning, one of those looks which illuminate the darkness  of the basest consciences. <I am astonished,> said he, <that,  thinking such things of M. Fouquet, you did not come to give me your counsels  thereupon.>

<Counsels upon what, sire?>

<Tell me, in the first place, clearly and precisely, what you think, M.  Colbert.>

<Upon what subject, sire?>

<Upon the conduct of M. Fouquet.>

<I think, sire, that M. Fouquet, not satisfied with attracting all the  money to himself, as M. Mazarin did, and by that means depriving your majesty  of one part of your power, still wishes to attract to himself all the friends  of easy life and pleasure—of what idlers call poetry, and politicians,  corruption. I think that, by holding the subjects of your majesty in pay, he  trespasses upon the royal prerogative, and cannot, if this continues so, be  long in placing your majesty among the weak and the obscure.>

<How would you qualify all these projects, M. Colbert?>

<The projects of M. Fouquet, sire?>

<Yes.>

<They are called crimes of lèse-majesté.>

<And what is done to criminals guilty of lèse-majesté?>

<They are arrested, tried, and punished.>

<You are quite certain that M. Fouquet has conceived the idea of the  crime you impute to him?>

<I can say more, sire; there is even a commencement of the execution of  it.>

<Well, then, I return to that which I was saying, M. Colbert.>

<And you were saying, sire?>

<Give me counsel.>

<Pardon me, sire; but in the first place, I have something to add.>

<Say—what?>

<An evident, palpable, material proof of treason.>

<And what is that?>

<I have just learnt that M. Fouquet is fortifying Belle-Isle.>

<Ah, indeed!>

<Yes, sire.>

<Are you sure?>

<Perfectly. Do you know, sire, what soldiers there are in  Belle-Isle?>

<No, ma foi! Do you?>

<I am ignorant, likewise, sire; I should therefore propose to your  majesty to send somebody to Belle-Isle?>

<Who?>

<Me, for instance.>

<And what would you do at Belle-Isle?>

<Inform myself whether, after the example of the ancient feudal lords, M.  Fouquet was battlementing his walls.>

<And with what purpose could he do that?>

<With the purpose of defending himself someday against his king.>

<But, if it be thus, M. Colbert,> said Louis, <we must  immediately do as you say; M. Fouquet must be arrested.>

<That is impossible.>

<I thought I had already told you, monsieur, that I suppressed that word  in my service.>

<The service of your majesty cannot prevent M. Fouquet from being  surintendant-general.>

<Well?>

<That, in consequence of holding that post, he has for him all the  parliament, as he has all the army by his largesses, literature by his favours,  and the noblesse by his presents.>

<That is to say, then, that I can do nothing against M. Fouquet?>

<Absolutely nothing,—at least at present, sire.>

<You are a sterile counsellor, M. Colbert.>

<Oh, no, sire; for I will not confine myself to pointing out the peril to  your majesty.>

<Come, then, where shall we begin to undermine this Colossus; let us  see;> and his majesty began to laugh bitterly.

<He has grown great by money; kill him by money, sire.>

<If I were to deprive him of his charge?>

<A bad means, sire.>

<The good—the good, then?>

<Ruin him, sire, that is the way.>

<But how?>

<Occasions will not be wanting; take advantage of all occasions.>

<Point them out to me.>

<Here is one at once. His royal highness Monsieur is about to be married;  his nuptials must be magnificent. That is a good occasion for your majesty to  demand a million of M. Fouquet. M. Fouquet, who pays twenty thousand livres  down when he need not pay more than five thousand, will easily find that  million when your majesty demands it.>

<That is all very well; I will demand it,> said Louis.

<If your majesty will sign the ordonnance I will have the money got  together myself.> And Colbert pushed a paper before the king, and  presented a pen to him.

At that moment the usher opened the door and announced monsieur le  surintendant. Louis turned pale. Colbert let the pen fall, and drew back from  the king, over whom he extended his black wings like an evil spirit. The  superintendent made his entrance like a man of the court, to whom a single  glance was sufficient to make him appreciate the situation. That situation was  not very encouraging for Fouquet, whatever might be his consciousness of  strength. The small black eye of Colbert, dilated by envy, and the limpid eye  of Louis XIV inflamed by anger, signalled some pressing danger. Courtiers are,  with regard to court rumours, like old soldiers, who distinguish through the  blasts of wind and bluster of leaves the sound of the distant steps of an armed  troop. They can, after having listened, tell pretty nearly how many men are  marching, how many arms resound, how many cannons roll. Fouquet had then only  to interrogate the silence which his arrival had produced; he found it big with  menacing revelations. The king allowed him time enough to advance as far as the  middle of the chamber. His adolescent modesty commanded this forbearance of the  moment. Fouquet boldly seized the opportunity.

<Sire,> said he, <I was impatient to see your majesty.>

<What for?> asked Louis.

<To announce some good news to you.>

Colbert, minus grandeur of person, less largeness of heart, resembled Fouquet  in many points. He had the same penetration, the same knowledge of men;  moreover, that great power of self-compression which gives to hypocrites time  to reflect, and gather themselves up to take a spring. He guessed that Fouquet  was going to meet the blow he was about to deal him. His eyes glittered  ominously.

<What news?> asked the king. Fouquet placed a roll of papers on the  table.

<Let your majesty have the goodness to cast your eyes over this  work,> said he. The king slowly unfolded the paper.

<Plans?> said he.

<Yes, sire.>

<And what are these plans?>

<A new fortification, sire.>

<Ah, ah!> said the king, <you amuse yourself with tactics and  strategies then, M. Fouquet?>

<I occupy myself with everything that may be useful to the reign of your  majesty,> replied Fouquet.

<Beautiful descriptions!> said the king, looking at the design.

<Your majesty comprehends, without doubt,> said Fouquet, bending  over the paper; <here is the circle of the walls, here are the forts,  there the advanced works.>

<And what do I see here, monsieur?>

<The sea.>

<The sea all round?>

<Yes, sire.>

<And what is, then, the name of this place of which you show me the  plan?>

<Sire, it is Belle-Île-en-Mer,> replied Fouquet with simplicity.

At this word, at this name, Colbert made so marked a movement, that the king  turned round to enforce the necessity for reserve. Fouquet did not appear to be  the least in the world concerned by the movement of Colbert, or the  king's signal.

<Monsieur,> continued Louis, <you have then fortified  Belle-Isle?>

<Yes, sire; and I have brought the plan and the accounts to your  majesty,> replied Fouquet; <I have expended sixteen hundred livres  in this operation.>

<What to do?> replied Louis, coldly, having taken the initiative  from a malicious look of the intendant.

<For an aim very easy to seize,> replied Fouquet. <Your  majesty was on cool terms with Great Britain.>

<Yes; but since the restoration of King Charles II. I have formed an  alliance with him.>

<A month since, sire, your majesty has truly said; but it is more than  six months since the fortifications of Belle-Isle were begun.>

<Then they have become useless.>

<Sire, fortifications are never useless. I fortified Belle-Isle against  MM. Monk and Lambert and all those London citizens who were playing at  soldiers. Belle-Isle will be ready fortified against the Dutch, against whom  either England or your majesty cannot fail to make war.>

The king was again silent, and looked askant at Colbert. <Belle-Isle, I  believe,> added Louis, <is yours, M. Fouquet?>

<No, sire.>

<Whose then?>

<Your majesty's.>

Colbert was seized with as much terror as if a gulf had opened beneath his  feet. Louis started with admiration, either at the genius or the devotion of  Fouquet.

<Explain yourself, monsieur,> said he.

<Nothing more easy, sire; Belle-Isle is one of my estates; I have  fortified it at my own expense. But as nothing in the world can oppose a  subject making an humble present to his king, I offer your majesty the  proprietorship of the estate, of which you will leave me the usufruct.  Belle-Isle, as a place of war, ought to be occupied by the king. Your majesty  will be able, henceforth, to keep a safe garrison there.>

Colbert felt almost sinking down upon the floor. To keep himself from falling,  he was obliged to hold by the columns of the wainscoting.

<This is a piece of great skill in the art of war that you have exhibited  here, monsieur,> said Louis.

<Sire, the initiative did not come from me,> replied Fouquet;  <many officers have inspired me with it. The plans themselves have been  made by one of the most distinguished engineers.>

<His name?>

<M. du Vallon.>

<M. du Vallon?> resumed Louis; <I do not know him. It is much  to be lamented, M. Colbert,> continued he, <that I do not know the  names of the men of talent who do honour to my reign.> And while saying  these words he turned towards Colbert. The latter felt himself crushed, the  sweat flowed from his brow, no word presented itself to his lips, he suffered  an inexpressible martyrdom. <You will recollect that name,> added  Louis XIV

Colbert bowed, but was paler than his ruffles of Flemish lace. Fouquet  continued:

<The masonries are of Roman concrete; the architects amalgamated it for  me after the best accounts of antiquity.>

<And the cannon?> asked Louis.

<Oh! sire, that concerns your majesty; it did not become me to place  cannon in my own house, unless your majesty had told me it was yours.>

Louis began to float, undetermined between the hatred which this so powerful  man inspired him with, and the pity he felt for the other, so cast down, who  seemed to him the counterfeit of the former. But the consciousness of his  kingly duty prevailed over the feelings of the man, and he stretched out his  finger to the paper.

<It must have cost you a great deal of money to carry these plans into  execution,> said he.

<I believe I had the honour of telling your majesty the amount.>

<Repeat it if you please, I have forgotten it.>

<Sixteen hundred thousand livres.>

<Sixteen hundred thousand livres! you are enormously rich,  monsieur.>

<It is your majesty who is rich, since Belle-Isle is yours.>

\begin{bwbigpic}
	[\basicscale] 
	{75rich}
	[<It is your majesty who is rich, since Belle-Isle is yours.>]
\end{bwbigpic}

<Yes, thank you; but however rich I may be, M. Fouquet\longdash> The  king stopped.

<Well, sire?> asked the superintendent.

<I foresee the moment when I shall want money.>

<You, sire? And at what moment then?>

<To-morrow, for example.>

<Will your majesty do me the honour to explain yourself?>

<My brother is going to marry the English Princess.>

<Well, sire?>

<Well, I ought to give the bride a reception worthy of the granddaughter  of Henry IV.>

<That is but just, sire.>

<Then I shall want money.>

<No doubt.>

<I shall want\longdash> Louis hesitated. The sum he was going to  demand was the same that he had been obliged to refuse Charles II. He turned  towards Colbert, that he might give the blow.

<I shall want, to-morrow\longdash> repeated he, looking at Colbert.

<A million,> said the latter, bluntly; delighted to take his  revenge.

Fouquet turned his back upon the intendant to listen to the king. He did not  turn round, but waited till the king repeated, or rather murmured, <A  million.>

<Oh! sire,> replied Fouquet disdainfully, <a million! what  will your majesty do with a million?>

<It appears to me, nevertheless\longdash> said Louis XIV.

<That is not more than is spent at the nuptials of one of the most petty  princes of Germany.>

<Monsieur!>

<Your majesty must have two millions at least. The horses alone would run  away with five hundred thousand livres. I shall have the honour of sending your  majesty sixteen hundred thousand livres this evening.>

<How,> said the king, <sixteen hundred thousand  livres?>

<Look, sire,> replied Fouquet, without even turning towards  Colbert, <I know that wants four hundred thousand livres of the two  millions. But this monsieur of l'intendance> (pointing over his  shoulder to Colbert, who if possible, became paler, behind him) <has in  his coffers nine hundred thousand livres of mine.>

The king turned round to look at Colbert.

<But\longdash> said the latter.

<Monsieur,> continued Fouquet, still speaking indirectly to  Colbert, <monsieur has received, a week ago, sixteen hundred thousand  livres; he has paid a hundred thousand livres to the guards, sixty-four  thousand livres to the hospitals, twenty-five thousand to the Swiss, an hundred  and thirty thousand for provisions, a thousand for arms, ten thousand for  accidental expenses; I do not err, then, in reckoning upon nine hundred  thousand livres that are left.> Then turning towards Colbert, like a  disdainful head of office towards his inferior, <Take care,  monsieur,> said he, <that those nine hundred thousand livres be  remitted to his majesty this evening, in gold.>

<But,> said the king, <that will make two millions five  hundred thousand livres.>

<Sire, the five hundred thousand livres over will serve as pocket money  for his royal highness. You understand, Monsieur Colbert, this evening before  eight o'clock.>

And with these words, bowing respectfully to the king, the superintendent made  his exit backwards, without honouring with a single look the envious man, whose  head he had just half shaved.

Colbert tore his ruffles to pieces in his rage, and bit his lips till they  bled.

Fouquet had not passed the door of the cabinet, when an usher pushing by him,  exclaimed: <A courier from Bretagne for his majesty.>

<M. d'Herblay was right,> murmured Fouquet, pulling out his  watch; <an hour and fifty-five minutes. It was quite true.>